% 本文件由 LaTeX 排版 Agent 根据有效 translation.json 自动生成。
% author=John Chrysostom; work_id=1917; title=致安提约基雅几位司铎书

\chapter{致安提约基雅几位司铎书}

% structure: p0001 source=h1 target=omitted reason=page-title-already-rendered-by-parent

以下这封信作为范例附上，出自数量众多的信件，展现了\Person{金口若望}写给亲密朋友时那种自然、近乎戏谑的风格和温暖亲切的语气。他现存的所有信件都是在流放期间写成的，因此内容多有重复，并具有很大的总体相似性。\par

\Emph{致\Person{卡斯图斯}、\Person{瓦莱里乌斯}、\Person{狄奥凡图斯}、\Person{西里亚库斯}——\Place{安提约基雅}的司铎们}\par

我不惊讶你们称我那封长信为短信。因为这正是恋爱中人的特点；他们不承认有满足这回事，也不接受饱足，反而从所爱之人那里得到的愈多，索求的也愈多。因此，即使你们收到的信比前一封大十倍，也逃不掉\Quote{简短}的称呼；事实上，它会被叫做短信，而且不仅被如此称呼，在你们眼中也真的如此。因此，我也同样对你们对我的爱所达到的尺度永不满足，而是总在寻求为你们的爱之饮添加分量，天天要求偿还那总在偿还却总在欠着的爱之债（因为经上记载：\BibleQuote{除了彼此相爱外，不可欠人什么}\BibleRef{罗 13:8}）。我确实不断大量领受我所求的，却从不以为已得全部。所以，不要停止偿付这笔好债务，它有着双重的喜乐。因为偿付者和收受者都同样享受，双方都因这偿付而一同变得富有；而在金钱上是不可能的，因为在那里，偿付者变穷，只有收受者变富。但在爱的盟约中，通常不是这样。因为偿付爱者并不因它被转移给收受者而变得匮乏，如金钱那样；相反，爱的偿付使偿付者比以前更富足。知道了这些之后，最受尊敬且虔诚的诸位先生，请你们不要停止向我不断显示这美好的情谊。因为虽然你们无需为此受我劝勉，但我极其渴望你们的爱，即使你们不需要，我也提醒你们，既要时常写信给我，也要告知我你们的健康状况。因为即使你们不需要任何人为此提醒你们，我也不会停止不断向你们寻求此事，这是我非常在意的事。我深知，由于季节、旅途艰难以及愿意为你们效劳的旅人稀少，这是一项艰巨的任务；但尽管如此，在如此多的困难中，在可能和可行的范围内，我们劝你们经常写信，并向你们的爱恳求这一恩惠。\par

\SourceRecord{Translated by W.R.W. Stephens.；Nicene and Post-Nicene Fathers, First Series；Vol. 9.；Edited by Philip Schaff.；Buffalo, NY: Christian Literature Publishing Co.,；1889.；<http://www.newadvent.org/fathers/1917.htm>.}
